
\section{引言}
\subsection{问题背景}

无线电频谱是现代信息社会的基础性资源，其秩序直接影响到民航导航、铁路调度、公众移动通信与应急指挥等关键业务的安全运行。然而随着无线设备的爆发式增长，人为或非人为的无线电干扰事件日益频发。非法设置的发射设备、故障的电子装置乃至部分工业设备的电磁泄漏，都可能在特定频段上形成持续干扰，轻则降低通信质量，重则导致关键业务链路中断。无线电干扰源的定位与清除由此成为无线电频谱管理中一项常态化且时效性极强的工作。

工程实践中，干扰源的定位通常分为初步监测、区域性排查与精准定位三个阶段。前两个阶段依托固定监测站、移动监测车或无人机等设备做大范围信号采集与分析，能够把干扰源的位置约束到一片相对较小的区域内；但受限于环境遮挡与低功率干扰源的传播特性，最终的精确定位仍需要技术人员携带便携式无线电方位监测测向机进入目标区域徒步排查，通过观察场强变化与方向特征并结合经验逐步逼近，最终确定干扰源位置并予以清除。

这一过程存在两个突出问题。其一是效率问题，技术人员徒步搜索一片半径以千米计的区域，耗时往往以小时甚至以天计，而干扰造成的业务中断却是以分钟计的经济损失。其二是安全问题，某些干扰源所处环境存在高压、辐射、易燃易爆或地形险峻等危险因素，人工排查本身构成对作业人员的安全威胁。因此，以机器人替代人工开展危险环境下的定位与清除作业，具有明确而迫切的需求。

\subsection{研究意义}

从社会意义看，搭载于机器狗的干扰源自动定位清除系统能够在危险环境中代替技术人员作业，既大幅缩短干扰清除的响应时间、降低业务中断造成的经济损失，又从根本上规避了人员进入危险区域的安全风险，对于提升无线电管理的自动化水平与应急保障能力具有直接的推动作用。从科研意义看，本题在数学上提出了一个具有相当普遍性的问题：在观测量只有带误差的角度、目标的个数与位置完全未知、且观测本身存在方向性盲区的条件下，如何设计一套既保证穷尽覆盖又追求时间最优的搜索与清除策略。这一问题在数学结构上涵盖了计算几何中的交会定位、统计估计中的误差传播与最优实验设计、组合优化中的覆盖问题与旅行商问题，以及不完全观测条件下的序贯决策，其建模思路可迁移到雷达组网定位、声源追踪、无人机区域搜索与传感器网络部署等广泛的工程场景。

\subsection{问题重述}

\textbf{问题1：}已知若干个检测点坐标及某干扰源关于这些检测点的示向度，给出采用交会定位法计算多边形定位区域直径即区域内任意两点之间距离最大值的算法，并回答以该定位区域直径为直径的圆能否覆盖此定位区域。

\textbf{问题2：}已知某全向干扰源在一个检测点处测得的示向度，给出第二个检测点的选择策略以获得关于该干扰源较好的定位效果，并给出第二个检测点的候选区域。

\textbf{问题3：}已知目标区域中的干扰源均为全向干扰源，总数在 $10\sim16$ 个之间但具体个数未知。现派出一只机器狗从原点出发，测向机初始频道为 1，对目标区域中的干扰源进行自动定位并清除。要求在确保所有干扰源被清除的前提下，完成任务的时间越短越好。制定机器狗自动搜索定位及清除的策略并设计相应算法，通过模拟器演练测试并给出正式测试结果。

\textbf{问题4：}目标区域中既有全向干扰源，又有定向干扰源，总数仍为 $10\sim16$ 个之间，定向干扰源的个数及定向方向也未知，其他条件与问题 3 相同。给出该条件下的检测与清除策略及算法，并通过模拟器演练测试检验其有效性。

\subsection{研究思路}

本文以"几何度量—布站优化—序贯决策—不完全观测下的鲁棒化"为主线展开。针对问题1，本文建立楔形半平面交模型，用凸多边形直径算法度量定位区域的尺度，并由 Thales 逆定理给出直径圆覆盖定位区域的充要条件，从而把"能否覆盖"这一模糊提问化归为可判定的几何命题。针对问题2，本文建立方位线扰动下的误差传播模型，先在垂直布站特例上导出闭式最优基线，再在可行域上以期望误差为目标做二维全局寻优，并以目标函数的次水平集给出候选区域。针对问题3，本文建立保证覆盖网模型以确定观测点的最少数量与最优布局，以旅行商下界刻画行程的极限，从而把"策略优劣"转化为"与下界的相对差距"。针对问题4，本文建立半圆盘可检测域模型，证明背向不可测定理并导出方位包围判据，据此求解零漏测冗余网的最小规模，并设计自适应补测策略以压低实际代价。这四个问题之间层层递进：问题1提供度量工具，问题2的结论直接决定问题3的逼近阈值，问题3的观测网又是问题4冗余网构造的基础。
